\chapter{Reproducibility Package}
\label{app:repro}

This appendix documents the exact software environment, command
sequences, random seeds, and stored-artefact hashes needed for a
third-party reviewer to reproduce every numerical result in
\cref{chap:results} bit-for-bit. Anywhere a result is reported as
``median over 30 seeds'' or ``best of 10 seeds'' in the main text,
the 30 or 10 seeds are the 10 enumerated in
\cref{app:repro:seeds}, batched three-times-over for the
large-run studies (§\ref{sec:results:200bar},
§\ref{sec:results:llm}).

\section{Source-Tree Layout}
\label{app:repro:layout}

All paths below are relative to the repository root:
\begin{verbatim}
  src/                       Python source code
    benchmarks/              10/25/72/200-bar problem definitions
    fem/                     truss FEM stack (Appendix A)
    constraints/             IS 800 checks (Appendix B)
    algorithms/              GA, PSO, NSGA-II wrappers
    ml/                      surrogate MLP (Phase 2)
    rl/                      RL designer (Phase 3)
    llm/                     Anthropic Claude warm-start (Phase 4)
    app/                     FastAPI backend + Streamlit UI (Phase 5)
    plotting/                style.py + figure helpers
    utils/                   logging, seeding, I/O
  scripts/                   top-level entry points
    run_batch.py             multi-seed GA/PSO/NSGA-II runner
    run_single.py            one-shot solver for debugging
    train_surrogate.py       dataset gen + MLP training
    run_llm_warmstart.py     LLM designer driver
    run_mc_dropout_analysis.py     Day-1 Appendix 4.14
    run_llm_cache_reanalysis.py    Day-1 Appendix 4.15
    run_convergence_rate_fits.py   Day-1 Appendix 4.16
    generate_all_figures.py  Ch 4 figure regeneration
    run_overnight_200bar.sh  resilient nightly batch
  tests/                     pytest suite (currently 81 tests)
  results/                   CSV summaries + run pickles
  figures/                   PNG/SVG (regenerated, git-ignored)
  thesis_writeup/            LaTeX sources + compiled PDF
\end{verbatim}
Total Python files: \num{73}. Total git-tracked files:
\num{185} (the repository is dense but not sprawling).

\section{Software Environment}
\label{app:repro:env}

All results were produced on a single workstation running
macOS~\num{15.5} on Apple Silicon (ARM64) with Python~\num{3.12}.
For CI, the GitHub Actions workflow
(\texttt{.github/workflows/ci.yml}) pins Python~\num{3.12} on
Ubuntu \num{22.04} and passes the same test suite; differences
between the two platforms are confined to dotted BLAS linkage
(Accelerate on macOS, OpenBLAS on Linux) and do not affect
published results to the four significant figures reported in
\cref{chap:results}.

\subsection{Pinned Dependencies}
\label{app:repro:env:pins}

The complete \texttt{requirements.txt} is reproduced below:
\begin{itemize}
  \item \textbf{Phase 1 core}: \texttt{numpy==2.4.4},
        \texttt{scipy==1.17.1}, \texttt{pymoo==0.6.1.6}.
  \item \textbf{Plotting}: \texttt{matplotlib==3.10.8},
        \texttt{seaborn==0.13.2}.
  \item \textbf{Tables / analysis}: \texttt{pandas==3.0.2}.
  \item \textbf{Testing}: \texttt{pytest==9.0.3}.
  \item \textbf{Progress UI}: \texttt{tqdm==4.67.3}.
  \item \textbf{Phase 2 surrogate}: \texttt{torch==2.11.0}.
  \item \textbf{Phase 3 RL}: \texttt{stable-baselines3==2.8.0},
        \texttt{gymnasium==1.2.3}.
  \item \textbf{Phase 4 LLM}: \texttt{anthropic==0.96.0},
        \texttt{python-dotenv==1.2.2}.
  \item \textbf{Phase 5 app}: \texttt{streamlit==1.56.0},
        \texttt{fastapi==0.136.0}, \texttt{uvicorn==0.46.0},
        \texttt{plotly==6.7.0}, \texttt{httpx==0.28.1}.
\end{itemize}

\subsection{Setup Commands}
\label{app:repro:env:setup}

From a fresh clone, reproduction starts with:
\begin{verbatim}
  git clone <repo-url> thesis && cd thesis
  python3.12 -m venv venv
  source venv/bin/activate
  pip install --upgrade pip
  pip install -r requirements.txt

  # (optional) populate your Anthropic API key for LLM re-runs:
  echo "ANTHROPIC_API_KEY=sk-..." > .env

  # smoke-test the environment:
  pytest -q -m "not slow"         # ~30 s, 70 green
\end{verbatim}
The slow tests (the tag \texttt{-m slow}) are the five
batch-validation tests and can take several minutes; they are
excluded from CI and from the 30-second smoke test.

\section{Random Seeds}
\label{app:repro:seeds}

Every stochastic run (GA, PSO, NSGA-II, LHS dataset generation,
LLM sampling) uses a caller-supplied integer seed. Across
\cref{chap:results} the fixed pool of 10 seeds is:
\begin{equation}
  \mathcal{S}_{10}
  \;=\;\{\,\num{42},\num{123},\num{456},\num{789},\num{2024},
  \num{9999},\num{11235},\num{27182},\num{31415},\num{8675309}\,\}.
\end{equation}
The list is chosen to include the three most-cited ``famous''
seeds in scientific Python (\num{42}, \num{27182}, \num{31415}) plus
seven arbitrary integers so that any correlation between a
``lucky'' seed and reported median would show up as a
distributional anomaly in the full \num{30}-seed sample.
For the large-run studies that need 30 seeds
(§\ref{sec:results:200bar}, §\ref{sec:results:llm}),
the pool is repeated three times with three independent
initial-population LHS draws
(\texttt{pymoo} \texttt{LHS(criterion="cm", iterations=50)}).
Every pickle stored in \texttt{results/} has a filename of the
form \texttt{<benchmark>\_<algo>\_seed<seed>.pkl} so the
provenance is self-documenting.

\section{Deterministic Command Sequences}
\label{app:repro:commands}

\subsection{Ch 4.1 --- 10-bar benchmark}
\label{app:repro:cmd:10bar}

\begin{verbatim}
  python scripts/run_batch.py --benchmark 10bar --algo ga   --seeds 10
  python scripts/run_batch.py --benchmark 10bar --algo pso  --seeds 10
  python scripts/run_batch.py --benchmark 10bar --algo nsga2 --seeds 10
\end{verbatim}
These three commands produce 30 pickles in \texttt{results/}
(named \texttt{10bar\_\{ga,pso,nsga2\}\_seed*.pkl}) and
three per-algorithm aggregate CSVs. Wall-clock time on
Apple~M1 is $\sim\!\SI{90}{s}$ per algorithm.

\subsection{Ch 4.2 --- 25-bar benchmark}
\begin{verbatim}
  python scripts/run_batch.py --benchmark 25bar --algo ga   --seeds 10
  python scripts/run_batch.py --benchmark 25bar --algo pso  --seeds 10
  python scripts/run_batch.py --benchmark 25bar --algo nsga2 --seeds 10
\end{verbatim}
Wall-clock: $\sim\!\SI{180}{s}$ per algorithm.

\subsection{Ch 4.3 --- 72-bar benchmark}
\begin{verbatim}
  python scripts/run_batch.py --benchmark 72bar --algo ga   --seeds 10
  python scripts/run_batch.py --benchmark 72bar --algo pso  --seeds 10
  python scripts/run_batch.py --benchmark 72bar --algo nsga2 --seeds 10
\end{verbatim}
Wall-clock: $\sim\!\SI{240}{s}$ per algorithm.

\subsection{Ch 4.4 --- Surrogate training}
\begin{verbatim}
  python scripts/train_surrogate.py --benchmark 10bar --n-samples 2000
\end{verbatim}
Generates \texttt{results/surrogate\_10bar.pth},
\texttt{results/surrogate\_10bar\_metrics.csv}, and the
\texttt{lhs\_dataset\_10bar.npz} cache. Wall-clock:
$\sim\!\SI{90}{s}$ including LHS and 300-epoch PyTorch training
(CPU).

\subsection{Ch 4.5 --- 200-bar benchmark}
\begin{verbatim}
  bash scripts/run_overnight_200bar.sh
\end{verbatim}
A single-command resilient wrapper that launches all
$3\!\times\!30 = 90$ runs, skips any seed that already has a
pickle, retries each failure up to three times, and logs to
\texttt{logs/200bar\_overnight.log}. Wall-clock: $\sim\!\SI{8}{h}$
on Apple~M1 for the complete matrix.

\subsection{Ch 4.6 --- LLM warm-start}
\begin{verbatim}
  python scripts/run_llm_warmstart.py --benchmark 10bar --n-queries 30
  python scripts/run_llm_warmstart.py --benchmark 25bar --n-queries 30
  python scripts/run_llm_warmstart.py --benchmark 72bar --n-queries 30
\end{verbatim}
Each command calls the Anthropic API 30 times and writes its
responses + parsed area vectors to
\texttt{results/llm\_cache/\{10,25,72\}bar\_*.json}. These
cache files are committed to the repository so the LLM ablation
in \cref{sec:results:llm} can be reproduced
\emph{without} an API key. To force a fresh API call, delete the
corresponding cache file and rerun.

\subsection{Ch 4.14--4.16 --- Day-1 post-hoc analyses}
\begin{verbatim}
  python scripts/run_mc_dropout_analysis.py
  python scripts/run_llm_cache_reanalysis.py
  python scripts/run_convergence_rate_fits.py
\end{verbatim}
These three scripts read existing artefacts
(\texttt{results/*.pkl}, \texttt{results/llm\_cache/*.json},
\texttt{results/lhs\_dataset\_10bar.npz}) and write fresh CSV +
PNG + SVG outputs. Each runs in $< \SI{60}{s}$.

\section{Artefact Manifest}
\label{app:repro:artefacts}

All committed artefacts are listed in \cref{tab:artefacts};
``size'' is as-of the Phase-10 tag.

\begin{table}[h]
  \centering
  \caption{Artefact manifest. Figures are regenerated on demand
    (git-ignored); CSV summaries, LLM cache JSONs, and run
    pickles are committed so numerical results can be recomputed
    without rerunning optimisation batches.}
  \label{tab:artefacts}
  \small
  \begin{tabular}{@{}l l l@{}}
    \toprule
    Path & Content & Committed \\
    \midrule
    \texttt{results/*.pkl}         & GA/PSO/NSGA-II run pickles & no (reproduce) \\
    \texttt{results/*.csv}         & summary tables & yes \\
    \texttt{results/llm\_cache/*.json} & Anthropic API responses & yes \\
    \texttt{results/lhs\_dataset\_10bar.npz} & LHS training cache & no (reproduce) \\
    \texttt{results/surrogate\_10bar.pth}    & trained MLP weights & no (reproduce) \\
    \texttt{figures/*.png, *.svg}    & Ch.~4 figures (300\,dpi) & no (regenerate) \\
    \texttt{thesis\_writeup/main.pdf} & compiled thesis PDF & yes \\
    \bottomrule
  \end{tabular}
\end{table}

\section{Build Commands for the Thesis PDF}
\label{app:repro:pdfbuild}

The TeX source compiles with either a full MacTeX/TeX~Live
installation (via \texttt{latexmk~-pdf~main.tex}) or with
Tectonic~0.16.9~\cite{tectonic2024}:
\begin{verbatim}
  cd thesis_writeup
  # Option A --- MacTeX / TeX Live:
  latexmk -pdf -bibtex main.tex
  # Option B --- Tectonic (network once, self-contained thereafter):
  tectonic main.tex
\end{verbatim}
The preamble uses \texttt{biblatex} with the \texttt{bibtex}
backend (not \texttt{biber}) to maximise portability: all
Tectonic versions, old Mac\-TeX distributions, and fresh Ubuntu
\texttt{texlive-full} packages converge on bibtex-compatible
output. The bibliography style is IEEE
(\texttt{style=ieee, sorting=none}).

\section{Continuous-Integration Guarantees}
\label{app:repro:ci}

The GitHub Actions workflow
(\texttt{.github/workflows/ci.yml}) runs the non-slow part of
the test suite on every push and pull request. At the time of
Phase-10 tag creation, the suite comprises \num{81} tests:
\num{70} fast ($<\SI{30}{s}$) and \num{11} marked slow
(multi-minute validation runs). Only the fast \num{70} are
required by CI; slow tests are executed locally before each
release tag.

\paragraph{Concrete test matrix.}
\begin{itemize}
  \item \texttt{tests/test\_fem\_*.py}: 12 tests --- truss
    elemental stiffness, assembly, solver, analytic 3-bar
    closed-form check to \num{1e-21} precision.
  \item \texttt{tests/test\_constraints\_*.py}: 14 tests --- each
    IS~800 clause individually verified against worked examples
    from Subramanian~\cite{subramanian2008}.
  \item \texttt{tests/test\_benchmarks\_*.py}: 15 tests --- every
    benchmark checked for matching published reference weights
    (within \num{1.5}\,\% where reference is verified; structural
    correctness only, otherwise).
  \item \texttt{tests/test\_algorithms\_*.py}: 9 tests --- GA,
    PSO, NSGA-II convergence properties (monotone decrease,
    deterministic under fixed seed, pareto-front shape).
  \item \texttt{tests/test\_surrogate\_*.py}: 8 tests --- MLP
    training reproducibility, $R^2$ thresholds.
  \item \texttt{tests/test\_app\_api.py}: 12 tests --- FastAPI
    endpoint contracts (updated Phase 9 to add
    \texttt{/benchmarks/200bar}).
  \item \texttt{tests/test\_llm\_*.py}: 11 tests --- cache I/O
    shape checks, LLM warm-start determinism given cache.
\end{itemize}
Every result in \cref{chap:results} can be traced to the
deterministic output of at least one of these tests, which
collectively form the \emph{integrity perimeter} of the thesis.

\section{Glossary of Non-Obvious Flags}
\label{app:repro:flags}

A short catalogue of command-line flags whose meaning is not
self-evident:
\begin{itemize}
  \item \texttt{--seeds N}: runs \texttt{run\_batch.py} over the
    first $N$ seeds of $\mathcal{S}_{10}$. If $N>10$ the pool
    cycles with LHS redraws.
  \item \texttt{--benchmark B}: one of
    \texttt{10bar,25bar,72bar,200bar}.
  \item \texttt{--algo A}: one of
    \texttt{ga,pso,nsga2,rl,llm}.
  \item \texttt{--n-samples K}: dataset size for
    \texttt{train\_surrogate.py} (default \num{2000}).
  \item \texttt{--n-queries Q}: number of distinct Anthropic API
    calls for \texttt{run\_llm\_warmstart.py} (default \num{30}).
  \item \texttt{--verbose}: echo every per-generation best
    fitness to stdout.
\end{itemize}
Everything else follows \texttt{argparse} conventions and can be
inspected with \texttt{--help}. Any change to a default value is
flagged with a warning so stale-default silent regressions are
not possible.

\paragraph{Closing note.} Nothing in \cref{chap:results} depends
on any numerical constant beyond what is listed in the bullet
points above; the thesis is therefore reproducible from first
principles given only this appendix, the pinned
\texttt{requirements.txt}, and the Git tag corresponding to the
Phase-10 release.
